\section{结论}
Muon 的精确方向可由谱范数单位球上的线性最小化解释，矩阵核范数自然出现在下降量中。通过动量误差递推，本文在同一框架内得到确定性 $O(T^{-1/2})$ 与固定批量随机 $O(T^{-1/4})$ 的平均梯度范数上界，并保留维度、噪声及初始目标差的依赖。这些结果说明在明确假设和参数调度下，归一化矩阵更新能够接近一阶平稳点，而非证明任意实用配置都会收敛到全局最优解。

在此基础上，本文补充了显式依赖谱范数近似误差的条件定理。统一的误差界 $\delta<1$ 使相对误差可吸收入下降项，但固定五步 NS 并不自动保证这一条件。实用混合版本的收敛仍需要与其 Nesterov、衰减和辅助参数组相匹配的联合分析。

新增实验中，病态矩阵分解没有显示 Muon 的普遍优势；CIFAR-10 同时给出三种子准确率、达标步数、H20 单步活跃时间与显存开销。共享服务器采用主动等待，墙钟反映本次资源策略而非独占设备吞吐量。更新谱只解释方向结构，不能直接推断表示特征坍缩；$K$ 消融展示本次观测到的权衡，不把 $K=5$ 预设为最优。后续可延长训练预算、增加随机种子，并研究动态步长、非光滑激活下的分析以及分布式通信成本。完整日志和公开版本使上述结论可以被复查。
